\setcounter{chapter}{4}
\setcounter{section}{0}

\chapter{Scenari futuri}
        Con l'avvento dell'Intelligenza Artificiale (IA) si sono sviluppate molte tecnologie prima ritenute impensabili, inoltre si sono semplificate molte procedure che in passato richiedevano operazioni complesse e in generale, ci sono molti nuovi servizi disponibili; molti software stanno investendo in agenti IA che possono aiutare l'utente a svolgere i suoi compiti, una sorta di assistente digitale. Nell'ambito dell'informatica è nata una nuova tipologia di programmazione il \textit{vibe coding}, ovvero descrivere l'applicazione che si vuole realizzare ad una IA e lasciare tutta la fase di scrittura del codice e creazione della logica ad essa, alcuni esempi possono essere Claude AI, Base44 e Nothing Essential Apps. Tuttavia il dibattito sulla vera utilità delle IA è sempre più animato: in molti hanno espresso il loro timore riguardo alla possibilità di perdere il posto di lavoro venendo soppiantati da IA. L'esempio più eclatante è dato dal movimento "Anti-AI" per quanto riguarda la generazione di immagini tramite IA Generative, ponendo in evidenza come questo sia ai limiti del plagio e furto, dato che i modelli che tutt'ora utilizziamo sono stati addestrati su una mole più che significativa di dati trovati online. 
    \par
        Per quanto riguardano le IA "tecniche" ovvero i vari \textit{Large Language Model} (\textit{LLM}) specializzati nella risoluzione di problemi matematici e problemi di programmazione, bisogna riconoscere che questi sono molto utili nel cercare soluzioni rapide a problemi complessi. Tuttavia ad oggi la loro accuratezza lascia ancora a desiderare, poiché molte delle soluzioni che propongono, soprattutto nei casi di risoluzione di codice, prevedono richiami a librerie inesistenti, errori logici, ripetizione di funzioni, operazioni inutili e molte altre operazioni non ottimizzate. \newline 
        Nel contesto di questa Tesi, la IA sarà sempre "tecnica" e può essere una salvezza come una condanna per la figura del Tecnico Informatico. 
        
    \par
        Prima di iniziare a descrivere quali possono essere le caratteristiche sia del Tecnico Informatico che le caratteristiche della IA, è necessario illustrare come del semplice codice può diventare una IA.\newline
    \par
        Quello che comunemente viene chiamato Intelligenza Artificiale è una branca dell'informatica che a sua volta si suddivide in molteplici materie, tra cui il \textit{Machine Learning}. Questa materia studia programmi, algoritmi e sistemi che possono "ragionare" come un essere umano: il sistema riceve degli input e, dopo un processo che viene chiamato addestramento, genera degli output calcolati da una serie di regole generate automaticamente. Esiste un'altra materia contenuta interna al \textit{Machine Learning}, ovvero il \textit{Deep Learning} dove i singoli algoritmi hanno reti neurali molto complesse che apprendono nuovi "metodi di ragionamento e adattamento" da una grande mole di dati eterogenei. \newline
        Prendendo spunto dalla natura, lo psicologo americano Frank Rosenblatt progetta il componente fondamentale di una rete neurale informatica: il \textit{perceptron}. Questa cellula ha una struttura simile a quella del neurone, molti ingressi, un nucleo e un singolo output; questo componente base elabora gli input forniti associandoli ad un peso che determina lui, nella fase dell'apprendimento. \newline
        Il perceptron ha un limite, può solamente risolvere problemi che sono linearmente scomponibili; per rimuovere questa limitazione, vengono inseriti più perceptron in serie creando così un \textit{perceptron multi-layer} (MLP). Espandendo questo ragionamento a più livelli si ottengono le attuali reti neurali, molti livelli che risolvono dei problemi scomposti dai neuroni precedenti. I livelli sono principalmente tre: il livello di input, ovvero i primi neuroni a ricevere i dati, il livello nascosto, la moltitudine di neuroni che eseguono le operazioni intermedie, e infine il livello di output, l'ultimo livello della rete neurale che fornisce poi il risultato finale.
    \par
        Le reti neurali possono funzionare secondo due tecniche principali, in base al loro scopo: la tecnica di regressione e la tecnica di classificazione. La tecnica di regressione ha lo scopo di fornire dei risultati predittivi in funzione dell'input; la tecnica di classificazione invece ha lo scopo di classificare il risultato in base agli input forniti. \newline
        Per poter scegliere la tecnica giusta bisogna prima definire il ciclo di vita di una rete neurale. Il ciclo è diviso in 5 fasi, la prima fase è la definizione della tecnica da usare in base alla tipologia di problema in cui si vuole applicare la rete neurale. La seconda fase riguarda invece la raccolta e la preparazione dei dati, normalizzandoli, in modo da non avere anomalie e migliorare l'efficienza dell'apprendimento; è importante che si scelga o si crei una mole di dati sufficientemente grande e sufficientemente eterogenea. La terza fase riguarda la configurazione del modello, scegliendo il numero di input, secondo anche i dati scelti precedentemente, e quali output deve fornire. Successivamente la quarta fase riguarda l'apprendimento, ovvero un processo iterativo dove la rete "impara" in base ai dati forniti; questa fase può essere supervisionata, con un'entità che ha l'output corretto e quindi può verificarne il risultato finale, oppure non supervisionato, lasciando la rete neurale agli output che crea; l'apprendimento non supervisionato è usato se non si dispongono dei risultati nel dataset fornito alla IA. Esiste anche l'apprendimento per rinforzo, ovvero che la rete neurale apprende in base a scenari reali o simulati, visualizzando le decisioni come ricompense o penalità, massimizzando poi le ricompense. \newline
        Infine, l'ultima fase del ciclo di vita riguarda la valutazione e successiva ottimizzazione della rete neurale; una volta terminata la fase di ottimizzazione, si può ripetere il ciclo ripartendo dall'addestramento della rete ottimizzata, proseguendo fino a che la rete ha raggiunto il grado di accuratezza richiesto\footnote{\cite{NeuralNetworks}}. 
    \par
        Un metodo alternativo alla creazione di una rete neurale è la neuroevoluzione; anche questo metodo prende spunto dalla natura con la mutazione casuale di geni, in questo contesto la mutazione casuale è un algoritmo di generazione randomica di valori numerici mentre i geni sono i pesi associati ad ogni input del singolo neurone. \newline
        La neuroevoluzione, oltre a gererare randomicamente i pesi tra più agenti, testa la rete neurale in un ambiente e poi ne viene valutato l'adattamento; in base alla valutazione, il neurone più adatto tramanda i suoi pesi a nuovi agenti generati e, questi pesi, subiscono una mutazione. \newline
        Questo tipo di creazione della rete neurale è utile se i risultati che si vogliono ottenere non hanno delle regole ben definite e chiare, tramite questo metodo, la rete impara facendo e arriverà ad un risultato ottimale\footnote{\cite{Neuroevolution}}. 
% ---------------------------------------------------------------------------------------
        \section{Le caratteristiche del Tecnico in loco}
            \par
                Come si può evincere da quanto fin qui esposto, la presenza in loco di un Tecnico Informatico è pressoché fondamentale per garantire la corretta erogazione del servizio, in questo caso universitario, situazione non possibile con una IA. \newline
                Un buon Tecnico sempre presente riuscirà sempre a portare a termine un intervento senza "esaurire i token" o dover rimandare a causa di un "sovraccarico del server" o per "indisponibilità del server". 
            \par
                Facendo riferimento agli interventi eseguiti, molti di questi potrebbero essere sì facilmente eseguibili dalle IA, ma bisogna prendere atto che queste non riuscirebbero ad agire in autonomia: per esperienza personale, la sola IA che riesce per ora a interagire con diversi applicativi è Google Gemini ma le app sono solamente quelle che Google mantiene aggiornate, oltre al fatto che molte delle funzionalità più avanzate sono bloccate da un abbonamento mensile. Molti dei modelli in commercio e d'uso anche quotidiano forniscono solamente una risposta, quasi sempre a favore di ciò che gli si chiede, fornendo una sorta di riassunto di diversa lunghezza con il suggerimento di possibili operazioni più o meno semplici.
            \par
                Bisogna sottolineare però la continua formazione da parte del Tecnico Informatico, in quanto molte delle tecnologie in uso oggi potrebbero non esserlo più tra qualche anno, un banale esempio in tal senso riguarda la Direttiva Europea che obbliga tutti i dispositivi mobili ad usare il cavo di ricarica USB tipo C, in modo da unificare i caricabatterie; entro aprile 2026 ogni portatile dovrà seguire questa normativa\footnote{\cite{DirettivaUE_2022/2380},\cite{EU-USB-C}}. 
% ---------------------------------------------------------------------------------------
        \section{Le caratteristiche delle IA}
            \par
                I vantaggi delle IA sono molteplici, tra cui sicuramente la rapidità nella risposta e la maggiore accuratezza nella risposta in relazione al prompt fornito, offrendo un'ampia gamma di molteplici soluzioni e indagando possibili cause. Hanno sicuramente uno spettro di competenze più ampio, attingendo dall'infinità di risorse che si trovano sul web, tanto che sotto numerosi punti di vista le IA possono essere considerate un'enciclopedia tascabile che permette di sapere tutto su qualsiasi tipologia di intervento. 
            \par
                Avendo molti modelli disponibili in commercio, si possono anche avere diverse "opinioni" da più modelli diversi, una sorta di consiglio decisionale, per ottenere il risultato più accurato possibile, rendendo l'intervento una sola messa in atto della risposta del consiglio delle IA.
            \par
                Tuttavia, le IA forniscono risposte che potrebbero contenere errori dovuti alla mole di dati su cui si sono addestrate e non è detto quindi che gli interventi proposti siano effettivamente corretti e adeguati alla vera problematica in oggetto. \newline
                In aggiunta, alle IA va individuato anche un eccessivo consumo inteso nell'immediato in relazione alla banda di rete dell'Università ma a lungo termine anche in riferimento alle conseguenze climatiche che si vengono a generare nel tenere accesi tutti i data center necessari, altamente energivori; per non parlare dell'attuale crisi di mercato per quanto riguarda il reperimento di componenti dei computer: nel giro di qualche mese, i prezzi delle memorie RAM si sono moltiplicati di oltre il 300\% del loro valore e i prezzi delle memorie NVME, in generale di memorie di massa ad alte prestazioni, sono saliti del circa 50\%. \newline
                Tuttavia basti questo accenno, in quanto analisi di mercato e discussione sulle tematiche ambientali non rientrano nel focus di questa tesi di Laurea. 
            \par
                Un problema più tangibile è la formazione di tutto il personale, nel caso in cui si voglia ad esempio adottare una IA come Tecnico Informatico: ogni membro del personale dovrebbe saper costruire un prompt degno di generare una risposta quanto più accurata al problema che si sta riscontrando e ogni membro del personale dovrebbe essere formato anche sulle basi dell'informatica, sia hardware sia software. 
% ---------------------------------------------------------------------------------------
        \section{Possibili soluzioni ibride}
            \par
                Nonostante questi sviluppi, la presenza di un Tecnico è ormai fondamentale e spesso accettata come un presupposto, in quanto basta solamente un primo contatto e la procedura d'intervento si avvierà subito, senza andare a scatenare una catena di tentativi e possibili errori che nuocerebbero alla continuità di servizio e all'immagine dell'Università. \newline 
                Riflettendo sul fatto che la maggior parte degli interventi di routine hanno riguardato collegamenti errati agli impianti audiovisivi risulta evidente che eseguire richieste alla IA per ogni singolo intervento sarebbe stato oneroso in termini di tempo, sia in riferimento alla durata dell'intervento stesso sia in relazione al tempo d'attesa della risposta da parte della IA. \newline
                Di seguito si possono trovare tre ipotesi di collaborazione a questa situazione, in cui il Tecnico si ritrova a convivere con uno strumento sempre più d'uso comune. 
% ---------------------------------------------------------------------------------------
            \subsection*{Collaborazione Pura 50 Tecnico - 50 IA}
                La prima ipotesi tende più a uno scenario ideale, ossia a una collaborazione pura tra il Tecnico e la IA. In questa situazione, la IA fornisce informazioni e suggerimenti su procedure e interventi mentre il Tecnico, oltre a eseguire l'intervento vero e proprio, pondera con cautela ciò che gli viene suggerito filtrando eventuali soluzioni multiple fornite dalla singola risposta della IA. Questo non solo renderebbe più rapidi gli interventi ma risolverebbe anche il problema delle piattaforme eterogenee presenti all'interno di quello che è un campus universitario.
% ---------------------------------------------------------------------------------------
            \subsection*{Collaborazione 70 Tecnico - 30 IA}
                La seconda ipotesi invece è più reale, dove il bilanciamento viene spostato verso l'esperienza del Tecnico. Non appena si verifica un problema che richiede un intervento, il Tecnico procede immediatamente a risolverlo in base alle sue conoscenze pregresse e informazioni che gli vengono fornite tramite il ticket di assistenza; solo nel caso in cui l'intervento si riveli più complesso del previsto, il Tecnico potrà consultare la IA formulando un prompt ben preciso, costruito sui tentativi già eseguiti, e provando ad applicare quanto gli viene proposto dalla riposta della IA. 
% ---------------------------------------------------------------------------------------
            \subsection*{Collaborazione 30 Tecnico - 70 IA}
                L'ultima ipotesi è la più improbabile: prevede un quasi completo affidarsi alla IA per la maggior parte degli interventi di ordinaria manutenzione, lasciando solamente gli interventi più delicati e complessi al Tecnico. Questo farebbe in modo da alleggerire sia l'aspetto delle responsabilità in caso di danni che l'aspetto della quantità di lavoro e quindi la suddivisione degli interventi in più giornate lavorative. \newline
                La IA, in questo caso, non potrà avere accesso ai file privati del personale, in quanto violerebbe i diritti sulla privacy. 
% ---------------------------------------------------------------------------------------
            \hfill
            \par
                Tali scenari e ipotesi di collaborazione possono funzionare solamente nel caso in cui la IA sia stata ben addestrata su una mole di dati eterogenea e soprattutto nell'ambito informatico e delle telecomunicazioni; questo ambito è in continuo e rapido sviluppo ma i parametri fondamentali e le tecniche d'intervento sono pressoché rimaste invariate. Parimenti gli interventi d'ordinaria manutenzione possono risultare spersonalizzanti e lasciarli in carico ad uno strumento automatizzato può migliorare la qualità complessiva degli interventi svolti. \newline
                Fare affidamento dunque a uno strumento così potente che permette, non solo semplificare i lavori di manutenzione, ma anche di sostituire l'intervento del Tecnico è sicuramente un passo nella giusta direzione; l'importante è che la figura del Tecnico non solo non venga completamente scalzata dalle nuove tecnologie ma anche che egli stesso sia in grado di filtrare con cura ciò che gli viene fornito come risposta e come "soluzione assoluta" dalla IA.